\section{Cross-Cutting Question: Architecture, or Prompting?}


The ablations in \S5.2 speak to \emph{what kind} of intervention could fix T1. Naming the hazard in the instruction and hiding it from the cameras both left the avoidance behavior among completing carries unchanged (hiding raises completion, an effect at the grasp) --- on its face neither a prompting nor a recognition gap, though the ablations are underpowered and the shield shows only that an \emph{oracle-fed} layer can restore clearance. A \textbf{powered promptability probe} points the same way: replacing the neutral instruction with an explicit safety command --- ``keep the blade away from the person,'' ``keep the cup away from the hot stove,'' ``stay clear of the person'' --- produced \textbf{no significant reduction in violations in any channel} at \emph{N} = 20--24 paired seeds per cell (T1 33 \% $\rightarrow$ 29 \%, McNemar \emph{p} = 1.0; T2 33 \% $\rightarrow$ 46 \%, \emph{p} = 0.58, carry-yaw range unchanged; T6 30 \% $\rightarrow$ 25 \%, \emph{p} = 1.0); the command shuffles \emph{which} episodes violate without moving the rate. Two cautions bound this: for T1 and T6 the unconditioned violation rate \emph{is} the completion rate, so the paired comparison can only detect a command that changes completion or yields a completing clear carry (none did; only T2 has headroom), and at \emph{N} = 24 pairs McNemar's exact test detects only a command that converts at least six of the eight violating pairs (33 \% $\rightarrow$ $\leq$ 8 \%). We therefore report a null for these commands at this power --- the same null SafeManip's three-level prompt ablation reports on GR00T \citep{ref30} --- and reserve a firmer statement for the non-ceiling design of \S7 (Fig. \ref{fig:fixability} summarizes the probe and the shield). The null is \textbf{not GR00T-specific}: on $\pi_{0.5}$/Franka an explicit spatial command naming the rendered hazard fails to reduce the plow-through (88 \% vs 94 \%, McNemar \emph{p} = 1.0, \emph{N} = 16) and degrades task success ($\approx$ 100 \% $\rightarrow$ 62 \%).

What the evidence supports is therefore a \textbf{hypothesis} we place at the center of the taxonomy: execution-phase failures are missing \emph{behavioral competences} --- absent from the imitation training distribution (\S2) rather than from the prompt or the percept --- addressable by architecture or an external layer (a CBF-style shield \citep{ref9}, a shielding layer \citep{ref10}, an orientation controller) rather than by better prompts. It is falsifiable, our evidence is consistent with it but not decisive, and testing it needs a \textbf{non-ceiling} design in which the ablations \emph{can} move the metric (\S7).

\textbf{Alternative views.} (i) \emph{This is collision avoidance under a new name.} The predicates are classical; the object of measurement --- a policy given no map, planner or filter, scored against human-referenced standards, with fixability ablations --- is not. (ii) \emph{A shield fixes it, so the policy need not learn it.} True for T1, and shown; repulsion does not fix T6 even at 0.80 m, the stop that fixes T6 never releases on a static person (\S5.4, \S5.7), neither addresses T2 or T5, and each needs its margin tuned per hazard --- the residual is what the policy must own and what sets the demand on the layer (\S1). (iii) \emph{The scenes are unavoidable by construction.} For T1 the shield is the witness that a clearing path exists and for T6 the stopped carry is; for T3a and T4 a violation-free reference trajectory is still missing, so we do not attribute those rates to the policy alone until it exists (\S7). (iv) \emph{The promptability null is a completion null.} For T1 and T6 it is; the claim rests on the non-saturated T2 and $\pi_{0.5}$ cells and on the unchanged corridor paths (Fig. \ref{fig:overlay}), and will stand or fall on the non-ceiling design.

